% Section 2: SV PDEs
\section{Review for Single Variable PDEs}
\label{sec:method}

Throughout, $\ddx u_j = (u_{j+1}-2u_j+u_{j-1})/\dx^2$ is the centered second difference and refinement means $\dx\to0$ with $\dt$ slaved to $\dx$.

\subsection{The scalar family and the optimal time step}

Consider the deterministic family
\begin{equation}\label{eq:family} u_t = D\,u_{xx} + \alpha\,u, \qquad D>0,\ \alpha\in\R,
\end{equation} whose stochastic motivation (the moment equations of \eqref{eq:spde}, with $\alpha = \frac{k(k-1)}2\sigma^2$) is deferred to \S\ref{sec:stochastic}. The pure diffusion case $\alpha=0$ is exactly Chu's setting \citep{chu2011}.

Apply forward Euler in time and the three-point stencil in space. Taylor expansion of both sides, followed by repeated differentiation of
\eqref{eq:family} to convert time derivatives of $u$ into space derivatives (the operators $D\partial_{xx}$ and $\alpha$ commute, giving the binomial identity $u_{t^k} = (D\partial_{xx}+\alpha)^k u$), yields the local truncation error organized by derivative type:
\begin{equation}\label{eq:lte}
\tau \;=\;
\underbrace{\Bigl(\tfrac{D^2\dt}2 - \tfrac{D\dx^2}{12}\Bigr)u_{xxxx}}_{\text{Group 1}}
\;+\;
\underbrace{\alpha D\,\dt\,u_{xx}}_{\text{Group 2}}
\;+\;
\underbrace{\tfrac{\alpha^2\dt}2\,u}_{\text{Group 3}}
\;+\;O(\dt^2,\dx^4).
\end{equation} Only Group 1 couples $\dt$ to $\dx$. Setting it to zero:

\begin{theorem}[Optimal time step {\citep{chu2011}}]\label{thm:ots}
For \eqref{eq:family} discretized by forward Euler / centered differences, the
unique step size annihilating the leading ($u_{xxxx}$) truncation term is
\begin{equation}\label{eq:dtopt2}
\boxed{\ \dto = \frac{\dx^2}{6D}\ }
\end{equation}
and no single step annihilates both the $u_{xxxx}$ and $u_{xxxxxx}$ truncation
terms, since the level-2 condition is $\dt = \dx^2/(\sqrt{60}\,D) \ne \dto$.
\end{theorem}
\noindent (Proof: one-line cancellation in \eqref{eq:lte}; uniqueness by the
level-2 computation. Full algebra in \texttt{sign-change.tex} \S1.)

\begin{proposition}[$\alpha$-independence]\label{prop:alpha}
$\dto$ is independent of the reaction parameter $\alpha$: the cancellation
operates on the diffusive group alone. For the stochastic hierarchy this makes
$\dto$ independent of the noise amplitude $\sigma$ and of the moment order $k$.
\end{proposition}

Following the derivation inspecting the truncation error, at $\dt=\dto$ the Group-1 term vanishes and the surviving residual is
\begin{equation}\label{eq:resid}
\tau\big|_{\dto} \;=\;
\frac{\alpha\dx^2}6\,u_{xx} + \frac{\alpha^2\dx^2}{12D}\,u \;+\; O(\dx^4).
\end{equation}

Then the truncation error $\tau$ as assembled in \eqref{eq:lte} is the \emph{forward} residual: $\text{exact} - \text{scheme} = \dt\,\tau + \cdots$. Cancelling it therefore requires \emph{adding} $\dt\,\tau$.

\begin{lemma}[Closed form correction term]The corrected factor at any given time step is,
\begin{equation}\label{eq:C1}
C^{(1)}_j = +\frac{\alpha\,\dx^4}{36D}\,\ddx u^n_j \;+\; \frac{\alpha^2\,\dx^4}{72D^2}\,u^n_j.
\end{equation}
\end{lemma}

\begin{theorem}[First-level accuracy]\label{thm:p1}

The corrected scheme
\begin{equation}\label{eq:scheme1} u^{n+1}_j = u^n_j + \dto\Bigl[D\,\ddx u^n_j + \alpha\,u^n_j\Bigr] + C^{(1)}_j
\end{equation} has step error $O(\dx^6)$ and global order $4$: the correction $C^{(1)}$ is itself $O(\dx^4)$, so its evaluation with the second-order operator $\ddx$ contributes only $O(\dx^6)$ to the step --- the three-point stencil is mathematically sufficient at $p=1$.
\end{theorem}
\noindent\emph{Status:} \tagV. Measured orders and errors, with the harness described in Appendix~\ref{app:repro}:

\begin{table}[h]
\centering
\caption{Scalar $p=1$ convergence, $u_t = Du_{xx} + \alpha u$ with
$D=0.1$, $\alpha\sigma^2=0.25$, Neumann BCs, $T=0.05$; $L^2$ error against the
exact cosine-series solution. Data: \texttt{code/implementation/results.md}
(generated by \texttt{convergence\_test.py}, 2026-07). All methods use the same
grid and step size $\dto$ at each level.}
\label{tab:scalar}
\begin{tabular}{@{}lcccc@{}}
\toprule
Method & $N=32$ & $N=64$ & $N=128$ & $N=256$ \\
\midrule
OTS-NIDC ($p=1$) & $8.54\text{e-}09$ & $5.26\text{e-}10$ & $3.27\text{e-}11$ & $2.04\text{e-}12$ \\
Forward Euler    & $1.23\text{e-}05$ & $3.02\text{e-}06$ & $7.52\text{e-}07$ & $1.88\text{e-}07$ \\
4th-order FD + FE& $1.52\text{e-}05$ & $3.79\text{e-}06$ & $9.45\text{e-}07$ & $2.36\text{e-}07$ \\
Reference expm   & $2.77\text{e-}05$ & $6.81\text{e-}06$ & $1.70\text{e-}06$ & $4.23\text{e-}07$ \\
\midrule
\multicolumn{5}{@{}l@{}}{\emph{Measured order (log--log slope, mean over levels):}}\\
OTS-NIDC & \multicolumn{4}{c}{$4.01$ \quad(FE / 4th-order FD / expm: $2.00$--$2.01$)}\\
\bottomrule
\end{tabular}
\end{table}

The 4th-order-FD row carries the paper's cleanest negative control: high-order spatial accuracy alone buys nothing when the step size is slaved to $\dx^2$ --- the $O(\dt)$ time error caps the whole scheme at second order. The boost must come from the time-step structure; that is the method's thesis, verified.


\textbf{The $p$-level tower and the stencil bound \todo{move this into an appendix, perhaps?}}

The general correction is generated by powers of the operator, with the sign of Remark~\ref{rem:sign}:
\begin{equation}\label{eq:tower} C^{(\le p)}_j \;=\; +\sum_{\ell=1}^{p} \frac{\dto^\ell}{(\ell+1)!}
\bigl(D\mathcal{D}_2 + \alpha\bigr)^{\ell} u^n_j,
\end{equation} where $\mathcal{D}_2$ is a spatial operator whose accuracy must be raised as $\ell$ grows. Spatial discretization error breaks the binomial symmetry of the time expansion, giving the stencil-scaling bound:

\begin{proposition}[Stencil bound; sketch]\label{prop:stencil}
Global order $2p+2$ requires corrections evaluated on at least a
$(2p+3)$-point stencil: $p=1$ needs the 3-point main stencil plus a 5-point
evaluation inside $C^{(1)}$ when full fourth-order spatial accuracy is demanded
of the correction itself; $p=2$ needs a 7-point stencil to represent the
uncancelled $u_{xxxxxx}$ residual.
\end{proposition}
\noindent\emph{Status:} \tagO for $p\ge2$ numerics (no harness yet; the bound's
derivation is a consistency argument in \texttt{OTS-NIDC.md} \S6). The $p=1$
row is \tagV via Table~\ref{tab:scalar}.

\subsection{Stability of Single Variable Systems}
\label{sec:stability}

\todo{phrase all of this as lemmas please god}

\subsubsection{Setting and criterion}
Set $u^n_j = G^n e^{ikj\dx}$. The three-point Laplacian has symbol
\[
\ddx u^n_j \;\to\; -\frac{4}{\dx^2}\,s\,u^n_j,\qquad s=\sin^2(k\dx/2)\in[0,1].
\]
Let $\mu = D\dt/\dx^2$, $\eta=\alpha\dx^2/D$ ($\alpha$ already contains
$\sigma^2$ for the moment equations), and $a=\alpha\dt=\mu\eta$.
The exact one-step factor is $E(s)=e^{-4\mu s+a}$, so $|E|\le e^{a}$: for
$\alpha>0$ the equation grows, and $|G|\le1$ is the wrong requirement.
We use the von Neumann condition with growth allowance,
\begin{equation}\label{eq:vN}
|G(s)|\le 1 + C\,\dt \quad\text{uniformly in } s\in[0,1],
\end{equation}
along refinement paths $\dt=\mu\dx^2/D$ with $\mu$ fixed (Lax--Richtmyer).
The sharper, natural comparison is $|G(s)|\le e^{a}$.

\subsubsection{Amplification factors}
We frame OTS-NIDC as a perturbation of the simple Finite-Euler method, so we begin by deriving the amplification factor $G_{\mathrm{FE}}$ and then transform it into $G_{\mathrm{NIDC}}$ as follows in,

\begin{align}
G_{\mathrm{FE}} &= 1-4\mu s+\mu\eta,\\
G_{\mathrm{NIDC}} &= 1-4\mu s+\mu\eta-\tfrac{\eta}{9}s+\tfrac{\eta^2}{72},
\end{align}
the latter from $C^{(1)}$ in \eqref{eq:C1}: $\frac{\alpha\dx^4}{36D}\cdot(-\frac{4s}{\dx^2})=-\frac{\eta}{9}s$ and $\frac{\alpha^2\dx^4}{72D^2}=\frac{\eta^2}{72}$. Both are affine in $s$ with negative slope, so the extremes over $s\in[0,1]$ are at $s=0$ (physical mode) and $s=1$ (grid mode).

\subsubsection{Stability at the optimal step}
At $\dt=\dto$, $\mu=1/6$, $\eta=6a$, and
\begin{equation}\label{eq:factor}
G_{\mathrm{NIDC}}(s)=\Bigl(1-\tfrac{2}{3}s\Bigr)(1+a)+\tfrac{a^2}{2}.
\end{equation}
Compare $E(s)=e^{-2s/3}e^{a}$: the correction supplies exactly the
$a^2/2$ term of $e^{a}$ at $s=0$, and $-\frac{2}{3}as$ is the cross term.
Since $1-\tfrac23 s\in[\tfrac13,1]$,
\[
\tfrac13(1+a)+\tfrac{a^2}{2}\;\le\;G_{\mathrm{NIDC}}(s)\;\le\;1+a+\tfrac{a^2}{2}\;\le\; e^{a}
\qquad(a\ge0).
\]
Hence $G_{\mathrm{NIDC}}>0$ for all $s$ (no oscillatory modes) and $|G_{\mathrm{NIDC}}|\le e^{\alpha\dt}$: \eqref{eq:vN} holds for every $\eta\ge0$, with no restriction on $\eta$. The offset $+\eta^2/72$ is therefore not destabilizing; it is what makes $G(0)$ match $e^{a}$ to $O(a^2)$. For plain FE at the same $\mu$, $G_{\mathrm{FE}}(s)=(1-\tfrac23 s)+a$, also inside $[\,\tfrac13+a,\,1+a\,]$.

\subsubsection{Grid-scale boundary for general $\mu$}
Requiring the lower side of \eqref{eq:vN} at $s=1$, i.e.\ $G(1)\ge-(1+\mu\eta)$:
\begin{align}
\mu_{\mathrm{FE},\max}&=\frac{2}{4-2\eta}=\frac{1}{2-\eta},\\
\mu_{\mathrm{NIDC},\max}&=\frac{2-\eta/9+\eta^2/72}{4-2\eta},
\end{align}
with ratio
\begin{equation}\label{eq:shrink2}
\frac{\mu_{\mathrm{NIDC},\max}}{\mu_{\mathrm{FE},\max}}
=1-\frac{\eta}{18}+\frac{\eta^2}{144}.
\end{equation}
(The stricter test $G(1)\ge-1$ gives denominators $4-\eta$ and the same ratio.) At $\eta=0.30$: $\mu_{\mathrm{FE},\max}=0.5882$, $\mu_{\mathrm{NIDC},\max}=0.5788$, ratio $0.98396$ ($-1.6\%$). The shrink is $O(\dx^2)$, and the effective NIDC CFL is $\dt\lesssim(1-\eta/18)\,\dx^2/(2D)$.

\subsubsection{Scope of the correction ($\dt\ne\dto$)}
$C^{(1)}$ is derived at $\dt=\dto$ and is independent of $\dt$. At $s=0$,
\[ G_{\mathrm{NIDC}}(0)-e^{a}\;\approx\;\frac{\eta^2}{72}\bigl(1-36\mu^2\bigr), \]
which is $\le0$ iff $\mu\ge 1/6$. For $\mu<1/6$ the excess is $\dt$-independent, so at fixed $\dx$ with $\dt\to0$ the accumulated growth $(1+\text{const})^{T/\dt}$ diverges.  Stability statements therefore refer to refinement paths $\mu$ fixed, where the excess is $O(\dt^2)$ and \eqref{eq:vN} holds; the $\dt/\dto\in\{0.5,\dots,2\}$ table is a sweep over such paths, not a fixed-$\dx$ study. The method is designed for $\mu=1/6$.

\subsubsection{Pair equations}
Replace $D\to D_i+D_j$, $\alpha\to\alpha_{ij}$; \eqref{eq:factor} holds verbatim with $\mu_{ij}=1/6$ at $\dto^{(ij)}$. For the coupled discrete system there is still no independent proof (open item \ref{open:stab}).


\begin{theorem}[Von Neumann stability]\label{prop:stability}
With $s=\sin^2(k\dx/2)$, $\mu = D\dt/\dx^2$, $\eta = \alpha\dx^2/D$, the amplification factors are $G_{\mathrm{FE}} = 1-4\mu s+\mu\eta$ and
\[ G_{\mathrm{NIDC}} = 1 - 4\mu s + \mu\eta - \tfrac{\eta}{9}s + \tfrac{\eta^2}{72}. \]
For the moment hierarchy $\eta \ge 0$ and the grid-scale stability bound tightens from $\mu \le 2/(4-\eta)$ to $\mu \le (2-\eta/9+\eta^2/72)/(4-\eta)$: a relative shrink of $1 - \eta/18 + O(\eta^2)$, i.e.\ $O(\dx^2)$, about $-1.6\%$ at the coarse $\eta=0.3$ test point. The optimal step $\mu = 1/6$ sits a factor of $3$ inside the CFL ceiling for both schemes.
\end{theorem}
\noindent\emph{Status:} \tagV (analytic derivation above + numerical check in
\texttt{notebooks/ai-src/stability/}; the earlier ``expands the CFL'' conclusion was an artifact of the pre-correction sign and has been fixed in the same release). The correction is stability-neutral to first order; it never threatens the optimal step. Note the honest asymmetry: this is a \emph{scalar} analysis; the coupled discrete schemes of \S\ref{sec:coupled} have no independent stability proof (open item \ref{open:stab}).

\subsection{Efficiency profile}

The method's cost argument is architectural, not just arithmetic: versus multi-stage Runge--Kutta, OTS-NIDC is single-pass (one read, one write of the state vector; no stage storage, no stage barriers); versus compact schemes, it is strictly local (no tridiagonal solves), so it maps onto GPUs and bandwidth-bound hardware without global communication. At $N=256$ the verified scalar runs show $\sim10^3$--$10^4\times$ lower error than forward Euler at identical grid and step \tagV --- order wins over constants at refinement.
